\documentclass[11pt]{article}
\usepackage[margin=1in]{geometry}
\usepackage{amsmath,amssymb,amsthm}
\usepackage[T1]{fontenc}
\usepackage{lmodern}
\usepackage{microtype}
\usepackage{bm}

\newtheorem{property}{Property}

\newcommand{\Ainv}{\bm{A}^{-1}}
\newcommand{\Binv}{\bm{B}^{-1}}
\newcommand{\A}{\bm{A}}
\newcommand{\B}{\bm{B}}
\newcommand{\I}{\bm{I}}
\newcommand{\tr}{\operatorname{tr}}
\newcommand{\rank}{\operatorname{rank}}

\title{\vspace{-1.5em}Matrix Demonstrations}
\author{Estimation I}
\date{}

\begin{document}
\maketitle
\vspace{-2.5em}

\noindent Throughout, $\A$ and $\B$ are conformable and nonsingular unless otherwise noted.

\section*{Inverse Properties}

\begin{property}[Scalar Inverse]
For $\lambda \neq 0$,
\[
(\lambda \A)^{-1} = \frac{1}{\lambda}\Ainv.
\]
\end{property}

\textbf{Demonstration:}
\begin{align*}
(\lambda\A)\left(\tfrac{1}{\lambda}\Ainv\right) &= \A\Ainv = \I, \\
\left(\tfrac{1}{\lambda}\Ainv\right)(\lambda\A) &= \Ainv\A = \I.
\end{align*}

\begin{property}[Inverse of a Product]
\[
(\A\B)^{-1} = \Binv\Ainv.
\]
\end{property}

\textbf{Demonstration:}
\begin{align*}
(\A\B)(\Binv\Ainv) &= \A(\B\Binv)\Ainv = \A\Ainv = \I, \\
(\Binv\Ainv)(\A\B) &= \Binv(\Ainv\A)\B = \Binv\B = \I.
\end{align*}

\begin{property}[Inverse Transpose]
\[
(\Ainv)' = (\A')^{-1}.
\]
\end{property}

\textbf{Demonstration:}
\begin{align*}
\A'(\Ainv)' &= (\Ainv\A)' = \I' = \I, \\
(\Ainv)'\A' &= (\A\Ainv)' = \I' = \I.
\end{align*}

\begin{property}[Sum Inverse]
\[
(\A+\B)^{-1} = \Ainv(\Ainv + \Binv)^{-1}\Binv.
\]
\end{property}

\textbf{Demonstration:} Factor $\A + \B$ by inserting identity matrices:
\begin{align*}
\A + \B &= \B\Binv\A + \B\Ainv\A
        = \B(\Binv + \Ainv)\A.
\end{align*}
Inverting both sides using Property~2:
\[
(\A+\B)^{-1} = \bigl(\B(\Ainv+\Binv)\A\bigr)^{-1}
             = \Ainv(\Ainv+\Binv)^{-1}\Binv.
\]

\begin{property}[Inverse Difference]
\[
\Ainv - (\A+\B)^{-1} = \Ainv(\Binv + \Ainv)^{-1}\Ainv.
\]
\end{property}

\textbf{Demonstration:} Factor $\A + \B$ differently:
\begin{align*}
\A + \B &= \A\Ainv\B + \A\Binv\B = \A(\Ainv + \Binv)\B.
\end{align*}
So $(\A+\B)^{-1} = \Binv(\Ainv+\Binv)^{-1}\Ainv$. Now write $\Binv = (\Binv + \Ainv) - \Ainv$ and expand:
\begin{align*}
(\A+\B)^{-1} &= \bigl[(\Binv+\Ainv) - \Ainv\bigr](\Binv+\Ainv)^{-1}\Ainv \\
              &= \I\cdot\Ainv - \Ainv(\Binv+\Ainv)^{-1}\Ainv \\
              &= \Ainv - \Ainv(\Binv+\Ainv)^{-1}\Ainv.
\end{align*}
Rearranging gives the result.

\section*{Structural Properties}

\begin{property}[Idempotent Matrices]
If $\bm{H}$ is a $k\times k$ idempotent matrix $(\bm{H}^2 = \bm{H})$, then $\bm{H}$ is either singular or the identity matrix, and $\tr(\bm{H}) = \rank(\bm{H})$.
\end{property}

\textbf{Demonstration (singular or identity):} If $\bm{H}$ is nonsingular, multiply $\bm{H}^2 = \bm{H}$ on the left by $\bm{H}^{-1}$:
\[
\bm{H}^{-1}\bm{H}^2 = \bm{H}^{-1}\bm{H} \implies \bm{H} = \I.
\]
So the only nonsingular idempotent matrix is $\I$.

\medskip

\textbf{Demonstration (trace equals rank):} Since $\bm{H}^2 = \bm{H}$, the eigenvalues of $\bm{H}$ must satisfy $\lambda^2 = \lambda$, so each eigenvalue is $0$ or $1$. The number of eigenvalues equal to $1$ is $\rank(\bm{H})$, and the trace is the sum of all eigenvalues, so $\tr(\bm{H}) = \rank(\bm{H})$.

\medskip

\textbf{Example:} The OLS projection matrix $\bm{P} = \bm{X}(\bm{X}'\bm{X})^{-1}\bm{X}'$ is idempotent:
\begin{align*}
\bm{P}^2 &= \bm{X}(\bm{X}'\bm{X})^{-1}\bm{X}'\bm{X}(\bm{X}'\bm{X})^{-1}\bm{X}' \\
          &= \bm{X}(\bm{X}'\bm{X})^{-1}\bm{X}' = \bm{P}.
\end{align*}
Note also that $\bm{P}' = \bm{X}((\bm{X}'\bm{X})^{-1})'\bm{X}' = \bm{X}(\bm{X}'\bm{X})^{-1}\bm{X}' = \bm{P}$, since $(\bm{X}'\bm{X})^{-1}$ is symmetric. So $\bm{P}$ is both idempotent and symmetric.

\begin{property}[Orthogonal Matrices]
If a $k\times k$ matrix $\bm{H}$ satisfies $\bm{H}'\bm{H} = \I$, then:
\begin{enumerate}
\item $\bm{H}$ is nonsingular with $\bm{H}^{-1} = \bm{H}'$,
\item $\bm{H}\bm{H}' = \I$,
\item $(\bm{H}')^{-1} = \bm{H}$.
\end{enumerate}
\end{property}

\textbf{Demonstration:}

(1) The equation $\bm{H}'\bm{H} = \I$ says that $\bm{H}'$ is a left inverse of $\bm{H}$. Taking determinants: $\det(\bm{H}')^2 = \det(\bm{H})^2 = 1$, so $\det(\bm{H}) \neq 0$ and $\bm{H}$ is nonsingular. Since the left inverse of a nonsingular matrix equals its inverse, $\bm{H}^{-1} = \bm{H}'$.

\medskip

(2) Multiply $\bm{H}'\bm{H} = \I$ on the left by $\bm{H}$ and on the right by $\bm{H}^{-1} = \bm{H}'$:
\[
\bm{H}\bm{H}'\bm{H}\bm{H}' = \bm{H}\bm{H}' \implies (\bm{H}\bm{H}')^2 = \bm{H}\bm{H}'.
\]
But more directly: $\bm{H}\bm{H}' = \bm{H}\bm{H}^{-1} = \I$.

\medskip

(3) From (1), $\bm{H}^{-1} = \bm{H}'$. By the inverse transpose property (Property~3), $(\bm{H}')^{-1} = (\bm{H}^{-1})^{-1} = \bm{H}$.

\section*{Matrix Inversion Lemmas}

The following identities are useful for updating inverses when a matrix is perturbed. Each is demonstrated by verifying that the proposed inverse, when multiplied by the original matrix, yields $\I$.

\begin{property}[Woodbury Matrix Identity]
For nonsingular $n\times n$ matrix $\A$, $n\times k$ matrix $\B$, nonsingular $k\times k$ matrix $\bm{C}$, and $k\times n$ matrix $\bm{D}$:
\[
(\A + \B\bm{C}\bm{D})^{-1} = \Ainv - \Ainv\B(\bm{C}^{-1} + \bm{D}\Ainv\B)^{-1}\bm{D}\Ainv.
\]
\end{property}

\textbf{Demonstration:} Let $\bm{R} = \Ainv - \Ainv\B(\bm{C}^{-1} + \bm{D}\Ainv\B)^{-1}\bm{D}\Ainv$. We verify $(\A + \B\bm{C}\bm{D})\bm{R} = \I$. Write $\bm{S} = \bm{C}^{-1} + \bm{D}\Ainv\B$ for brevity.
\begin{align*}
(\A + \B\bm{C}\bm{D})\bm{R}
&= \I - \B\bm{S}^{-1}\bm{D}\Ainv + \B\bm{C}\bm{D}\Ainv - \B\bm{C}\bm{D}\Ainv\B\bm{S}^{-1}\bm{D}\Ainv \\
&= \I + \B\bigl[\bm{C}\bm{D}\Ainv - \bm{S}^{-1}\bm{D}\Ainv - \bm{C}\bm{D}\Ainv\B\bm{S}^{-1}\bm{D}\Ainv\bigr] \\
&= \I + \B\bigl[\bm{C}\bm{D}\Ainv - (\I + \bm{C}\bm{D}\Ainv\B)\bm{S}^{-1}\bm{D}\Ainv\bigr].
\end{align*}
Now $\I + \bm{C}\bm{D}\Ainv\B = \bm{C}(\bm{C}^{-1} + \bm{D}\Ainv\B) = \bm{C}\bm{S}$, so:
\begin{align*}
&= \I + \B\bigl[\bm{C}\bm{D}\Ainv - \bm{C}\bm{S}\bm{S}^{-1}\bm{D}\Ainv\bigr] \\
&= \I + \B\bigl[\bm{C}\bm{D}\Ainv - \bm{C}\bm{D}\Ainv\bigr] = \I. \qed
\end{align*}

\begin{property}[Sherman--Morrison Formula]
For nonsingular $\A$ and column vector $\bm{b}$, provided $1 + \bm{b}'\Ainv\bm{b} \neq 0$:
\[
(\A + \bm{b}\bm{b}')^{-1} = \Ainv - \frac{\Ainv\bm{b}\bm{b}'\Ainv}{1 + \bm{b}'\Ainv\bm{b}}.
\]
\end{property}

\textbf{Demonstration:} This is the Woodbury identity with $\B = \bm{b}$, $\bm{C} = 1$, $\bm{D} = \bm{b}'$:
\[
(\A + \bm{b}\cdot 1\cdot\bm{b}')^{-1}
= \Ainv - \Ainv\bm{b}(1 + \bm{b}'\Ainv\bm{b})^{-1}\bm{b}'\Ainv
= \Ainv - \frac{\Ainv\bm{b}\bm{b}'\Ainv}{1 + \bm{b}'\Ainv\bm{b}}.
\]
The subtraction variant follows by replacing $\bm{b}$ with $i\bm{b}$ (where $i^2 = -1$), giving:
\[
(\A - \bm{b}\bm{b}')^{-1} = \Ainv + \frac{\Ainv\bm{b}\bm{b}'\Ainv}{1 - \bm{b}'\Ainv\bm{b}},
\]
provided $1 - \bm{b}'\Ainv\bm{b} \neq 0$.

\end{document}
